\documentclass[11pt]{article}

% ---------- Packages ----------
\usepackage[margin=1in]{geometry}
\usepackage[T1]{fontenc}
\usepackage{lmodern}
\usepackage{microtype}
\usepackage{amsmath, amssymb, amsthm, mathtools}
\usepackage{bm}
\usepackage{bbm}
\usepackage{physics}
\usepackage{enumitem}
\usepackage{booktabs}
\usepackage{siunitx}
\usepackage{graphicx}
\usepackage{caption}
\usepackage[hidelinks]{hyperref}
\usepackage[nameinlink,capitalize]{cleveref}

% ---------- Theorem-like environments ----------
\newtheorem{theorem}{Theorem}
\newtheorem{proposition}[theorem]{Proposition}
\newtheorem{corollary}[theorem]{Corollary}
\theoremstyle{definition}
\newtheorem{definition}[theorem]{Definition}
\theoremstyle{remark}
\newtheorem{remark}[theorem]{Remark}

% ---------------------------
% Shortcuts & notation
% ---------------------------
\DeclareMathOperator{spec}{spec}
\newcommand{\RR}{\mathbb{R}}
\newcommand{\CC}{\mathbb{C}}
\newcommand{\PP}{\mathbb{P}}
\newcommand{\opnorm}[1]{\left\|#1\right\|}
\newcommand{\HS}{\mathrm{HS}}
\newcommand{\N}{\mathbb{N}}
\newcommand{\R}{\mathbb{R}}
\newcommand{\C}{\mathbb{C}}
\newcommand{\Pset}{\mathbb{P}}
\newcommand{\Hcal}{\mathcal{H}}
\newcommand{\absval}[1]{\left\lvert #1 \right\rvert}
\newcommand{\1}{\mathbbm{1}}
\newcommand{\gap}{\operatorname{gap}}
\newcommand{\LB}{\mathrm{LB}}
\newcommand{\UB}{\mathrm{UB}}


\title{\vspace{-1em}DRMM $\times$ CSC:\\
Certified Spectral Control of Prime-Scripted Operators (Finite-$P$)}
\author{Dr. Ryan Van Gelder, Luis Morató de Dalmases, \\ Tyler Van Osdol}
\date{\today}

\begin{document}
\maketitle
\vspace{-2em}

% =========================
% Motivation (Front-matter)
% =========================
\section*{Motivation \& Long-Term Vision}
We aim to develop a \emph{certifiable control theory for prime-scripted spectral operators} that (i) produces verifiable, finite-dimensional results now; (ii) scales to structurally faithful surrogates of operators from analytic number theory (Hecke/trace, automorphic Laplacians, moonshine blocks); and (iii) yields \emph{lawfulness-aware} design principles---prime gating and recursion budgets---that survive model variation. Success is \emph{not} a theorem like RH; it is a \emph{reliable pipeline} that (1) constructs operators with arithmetic provenance, (2) moves specific gaps/slopes under explicit budgets, and (3) quantifies when DRMM-style constraints do or do not confer real robustness beyond ordinary regularization.

% =========================
% Abstract
% =========================
\begin{abstract}
We present a finite-$P$, certified control scheme for prime-scripted operators that constrains the \emph{Certified Spectral Control} (CSC) design by \emph{DRMM} (Dynamic Recursive Meta-Mathematics) lawfulness budgets and introduces a small-gain DRMM recursion. The plant is $\mathcal U_P(\omega)=X_P + C(\omega;w)$ on $\ell^2(\PP)$, where $X_P$ encodes arithmetic symmetry and $C(\omega;w)$ is a Bohr--prime multiplier driven by weights $w$. We provide convex design surrogates (gap lower bound; slope upper bound), a pinning \& slope-control theorem with explicit constants, a lawful-controller program, a minimal $P{=}5$ sanity check with verification \& falsifiability, and a bounded recursion with a small-gain certificate. We claim \emph{finite-$P$} certified control and a clear scaling path---not asymptotic problem resolution.
\end{abstract}

% =========================
% 1. Plant & sector
% =========================
\section{Plant, Sector, and Design Objectives}
\paragraph{Indexing and space.}
Let $\PP=\{p_1,\dots,p_P\}$ be the first $P$ primes and work on $\Hcal=\ell^2(\PP)$.

\paragraph{Plant.}
The controlled operator is
\begin{equation}
\mathcal U_P(\omega)=X_P + C(\omega;w),\qquad
C(\omega;w)=\sum_{p\in\PP} w_p\, B_p(\omega),
\end{equation}
with $X_P$ self-adjoint, each channel $B_p(\omega)$ self-adjoint, and $w\in\R^P$ the controller. Assume uniform bounds $\norm{B_p(\omega)}\le b_p$ and $\norm{\partial_\omega B_p(\omega)}\le b'_p$, so
\begin{equation}
\norm{C(\omega;w)}\le \sum_p |w_p|\,b_p =:\norm{w}_{1,b},\quad
\norm{\partial_\omega C(\omega;w)}\le \sum_p |w_p|\,b'_p =:\norm{w}_{1,b'}.
\end{equation}

\subsection{Lawful sector (DRMM) --- constraints with motivation}\label{sec:lawful}
We constrain design and evolution to a \emph{lawful} subset $\Hcal_{\mathrm{lawful}}\subseteq\Hcal$ via three budgets.

\paragraph{(a) Prime gating ($\Lambda_m$): scale separation by prime size.}
For $\alpha$ in a small grid $\mathcal A$ (e.g., $\{0.5,0.7,1.0,1.3\}$) and $\kappa\in\R$, penalize deviations from a power-law template
\begin{equation}
R_{\Lambda}^{(\alpha)}(w,\kappa)=\sum_{p\in\PP}\big|\,w_p-\kappa p^{-\alpha}\,\big| \le \tau.
\end{equation}
We keep convexity by optimizing $(w,\kappa)$ for each fixed $\alpha\in\mathcal A$ and selecting $\alpha$ by validation.

\paragraph{(b) $\Xi$-budget: memory/recursion gain.}
A recursion operator $\Xi$ (introduced in \cref{sec:recursion}) satisfies $\norm{\Xi}\le \beta$. The pilot uses $\beta=0$; later, $\beta>0$ is certified by a small-gain margin.

\paragraph{(c) Commutator stability ($M$): channel geometry control.}
Bound non-abelian channel interference by
\begin{equation}
\sum_{p<q}\norm{[B_p,B_q]} \le \gamma,
\end{equation}
implemented via epigraph variables $\theta_{pq}\ge \norm{[B_p,B_q]}$ with $\sum_{p<q}\theta_{pq}\le\gamma$ (convex).

\paragraph{Design surrogates.}
Let $\lambda_1(\omega)\le\cdots\le\lambda_P(\omega)$ be the eigenvalues of $\mathcal U_P(\omega)$. If $X_P$ has a simple eigenvalue $\lambda_k$ separated by $\delta_S>0$, Weyl's inequality yields the \emph{gap lower bound}
\begin{equation}
\underline{\gap}_k(w):=\delta_S - 2\norm{w}_{1,b},
\end{equation}
and differentiability of $B_p$ yields the \emph{slope upper bound}
\begin{equation}
\Big|\frac{d\lambda_k}{d\omega}\Big| \le \norm{w}_{1,b'}.
\end{equation}

% =========================
% 2. Pilot theorem
% =========================
\section{Pilot Theorem: Pinning, Gap, Slope, and Subspace Control}

\paragraph{Assumptions A.}
(i) $X_P$ self-adjoint with a simple $\lambda_k$ such that $\min(\lambda_{k+1}-\lambda_k,\lambda_k-\lambda_{k-1})\ge \delta_S>0$.  
(ii) $\norm{C(\omega;w)}\le r<\delta_S/2$ for all $\omega$.  
(iii) $\partial_\omega C$ exists with $\norm{\partial_\omega C}\le s$.

\begin{theorem}[Pinning \& Slope Control]
Under Assumptions A:
\begin{enumerate}[label=(\roman*)]
\item \textbf{Pinning \& gap:} For all $\omega$, there is a simple eigenvalue $\lambda_k(\omega)$ with
$|\lambda_k(\omega)-\lambda_k|\le r$ and $\gap_k(\omega)\ge \delta_S-2r$.
\item \textbf{Slope:} With a smooth eigen-branch,
$\absval{d\lambda_k/d\omega}\le s$.
\item \textbf{Subspace stability:} If $\mathsf P_k$ and $\mathsf P_k(\omega)$ denote the rank-1 spectral projectors at $\lambda_k$ and $\lambda_k(\omega)$, then (Davis--Kahan)
$\norm{\mathsf P_k(\omega)-\mathsf P_k}\le \dfrac{r}{\delta_S-r}$.
\end{enumerate}
\end{theorem}

\begin{corollary}[Certified design via convex surrogates]
If $\norm{w}_{1,b}\le r<\delta_S/2$ and $\norm{w}_{1,b'}\le s$, any optimizer of the controller program in \cref{sec:controller} certifies the above bounds; a single eigentracking sweep verifies realized values.
\end{corollary}

% =========================
% 3. Controller
% =========================
\section{Controller: Convex Program with Lawfulness Budgets}\label{sec:controller}
Maximizing $\underline{\gap}_k(w)$ is equivalent to \emph{minimizing} $\norm{w}_{1,b}$ subject to budgets:
\begin{align}
\min_{w,\kappa}\quad & \norm{w}_{1,b} \nonumber\\
\text{s.t.}\quad
& \norm{w}_{1,b}\le r,\quad \norm{w}_{1,b'}\le s, \label{eq:prog}\\
& \exists\,\alpha\in\mathcal A:\ R_{\Lambda}^{(\alpha)}(w,\kappa)\le \tau, \nonumber\\
& \norm{\Xi}\le \beta,\quad \sum_{p<q}\theta_{pq}\le \gamma,\ \ \theta_{pq}\ge \norm{[B_p,B_q]},\nonumber\\
& w\in\R^P,\ \kappa\in\R.\nonumber
\end{align}
\textbf{Solver:} projected/proximal gradient for each fixed $\alpha\in\mathcal A$; select $\alpha$ by validation (e.g., realized gap at matched budgets). Warm start $w_p=\kappa p^{-\alpha}$.

% =========================
% 4. Verification
% =========================
\section{Verification \& Falsifiability}
\paragraph{Eigensolver.} One Lanczos sweep across a grid $\Omega=\{\omega_j\}$ around the design point; use Woodbury updates when channels are low rank.

\paragraph{Certified vs realized table.}
\begin{table}[h]
\centering
\caption{Certified vs.\ realized metrics at $\omega_j\in\Omega$.}
\begin{tabular}{@{}lcccc@{}}
\toprule
$\omega_j$ & $\underline{\gap}_k(w)$ (cert) & $\gap_k(\omega_j)$ (real) & $s$ (cap) & $\abs{\Delta\lambda_k/\Delta\omega}$ (real) \\
\midrule
$\omega_1$ & \dots & \dots & \dots & \dots \\
$\omega_2$ & \dots & \dots & \dots & \dots \\
\bottomrule
\end{tabular}
\end{table}

\paragraph{Null controls.} Generate $N$ null controllers $w^{\text{null}}$ by (i) shuffling prime lags in $B_p$, or (ii) randomizing signs, while matching $\norm{w}_{1,b}$ and $R_\Lambda$. Plot a histogram of gap improvements and report a $p$-value for our controller relative to nulls.

% =========================
% 5. Minimal experiment & scaling
% =========================
\section{Zeroth-Order Sanity Check (P=5) \& Scaling Plan}
\paragraph{Index set.} $\PP=\{2,3,5,7,11\}$.

\subsection{Base-operator menus (choose one per run)}
\paragraph{(H) Hecke-trace truncation (Automorphic-DRMM).}
$X_P=\sum_{q\le Q}\rho_q\,T_q$ restricted to $\ell^2(\PP)$ with symmetrized Hecke-like $T_q$.
\emph{Why:} arithmetic symmetry; $\delta_S$ measurable and tunable via $\rho_q$.

\paragraph{(E) Explicit-formula surrogate (pair-correlation toy).}
$(X_P)_{ij}=\psi(\log p_i)+\psi(\log p_j)-\psi(\log (p_ip_j))-\mu$ with smooth $\psi$.
\emph{Why:} mimics explicit-formula interactions while remaining bounded/symmetric.

\paragraph{(M) Moonshine block (modular moment matrix).}
$X_P=\Pi^\top M(\tau_0)\Pi$, where $M(\tau_0)$ is a small covariance of modular coefficients; $\Pi$ projects to prime indices.
\emph{Why:} direct modular provenance.

\paragraph{House rule.} Justify the chosen $X_P$ in two sentences and publish $\delta_S$.

\subsection{Channels and bounds}
Rank-1 Bohr--prime multipliers
\begin{equation}
B_p(\omega)=u_p(\omega)u_p(\omega)^\top,\qquad
u_p(\omega)_i=\cos\!\big(\omega\,\log(pp_i+\ell)\big),
\end{equation}
so $\norm{B_p}\le 1$ and $\norm{\partial_\omega B_p}\le b'_p$ (explicit from $u_p$).

\subsection{Protocol \& outputs}
(i) Solve \eqref{eq:prog} with published $(r,s,\tau,\beta,\gamma,\alpha)$.  
(ii) Verify by Lanczos across $\Omega$.  
(iii) Report: confidence-band plot and the certified vs.\ realized table; include null-control histogram and $p$-value.

\subsection{Scaling targets \& complexity}\label{sec:scaling}
Targets: $P=5$ (sanity), $P=23$ ($\sim$9$\times$), $P=97$ ($\sim$19$\times$).  
Controller: proximal methods with $O(P)$ primitives/step (empirically tens of iterations).  
Eigensweep: Lanczos $O(P\cdot m)$ matvecs (tens--low hundreds of Krylov steps for $P\le 100$).  
\emph{Success at scale:} certified bands remain binding (non-vacuous) and realized improvements beat nulls for $P\in\{23,97\}$.

% =========================
% 6. Recursion (DRMM)
% =========================
\section{Pilot-Plus: Bounded DRMM Recursion with Small-Gain Certificate}\label{sec:recursion}
Introduce a recursion on $T_t$:
\begin{equation}
T_{t+1}=T_t+\varepsilon\sum_{p\in\PP}\Lambda_m\,\Xi\,M_p(T_t),\qquad T_0=X_P+C(\omega;w),
\end{equation}
with $\norm{\Xi}\le \beta$ and mixer $M_p$ chosen to reduce channel conflict.

\paragraph{Default mixer.}
Let $u_p$ denote the channel vector underlying $B_p=u_pu_p^\top$. Define a Householder reflector
$U_p:=I-2\,\dfrac{u_pu_p^\top}{\norm{u_p}^2}$ and set
\begin{equation}
M_p(T):=U_pTU_p^\top - T,\qquad \norm{M_p(T)}\le 2\norm{T}.
\end{equation}
With $m_p:=2\norm{T_t}$, define the one-step recursion gain
\begin{equation}
\eta:=\varepsilon\sum_{p\in\PP}\Lambda_m\,\beta\,m_p.
\end{equation}

\begin{proposition}[Spectral drift control]
Suppose \cref{eq:prog} yields margin $\delta_S-2r>0$ and $\eta<\delta_S/8$. Then
\[
\norm{\mathsf P_k(T_{t+1})-\mathsf P_k(T_t)}\le \frac{\eta}{\delta_S-2r-\eta},\qquad
\gap_k(T_{t+1})\ge \delta_S-2r-2\eta.
\]
Over $T$ steps, pinning persists provided $T\eta<\delta_S/4$ (triangle inequality/telescoping).
\end{proposition}

\paragraph{Reporting.} Publish $(\varepsilon,\beta,\Lambda_m)$, per-step projector drift, cumulative drift over $T$, and realized gap floors vs.\ certified margins.

% =========================
% 7. Reproducibility & hyperparams
% =========================
\section{Reproducibility Checklist \& Hyperparameter Protocol}
\paragraph{Checklist.} Publish $P$, choice of $X_P$, channel definitions, and $(\phi,\mu)$ or $(\rho_q,Q)$ as applicable; list $(r,s,\tau,\beta,\gamma,\alpha)$; fix seeds; release solver and verification scripts.

\subsection{Setting $(r,s,\tau,\beta,\gamma)$ from data}
\begin{enumerate}[leftmargin=2em]
\item $r=\theta\,\delta_S$ with $\theta\in[0.2,0.4]$; publish $(\delta_S,\theta)$.
\item $s$: compute $b'_p$ analytically; enforce $\norm{w}_{1,b'}\le s$ with target $s\le 0.1\,\delta_S$.
\item $\tau$: for each $\alpha\in\mathcal A$, fit $\kappa$ by least-absolute-deviation; set $\tau=\eta\sum_p p^{-\alpha}$ with $\eta\in\{0.02,0.05\}$.
\item $\beta$: start at $0$; for recursion, choose $\beta$ so that $\eta<\delta_S/8$.
\item $\gamma$: measure $\sum_{p<q}\norm{[B_p,B_q]}$ on the uncontrolled system; set $\gamma$ to its $25$th percentile (conservative geometry).
\end{enumerate}

% =========================
% 8. Variants
% =========================
\section{Extension Tracks (Problem-Themed Variants)}
\textbf{Automorphic-DRMM:} $X_P=\sum_{q\le Q}\rho_qT_q$; test with/without $(\Lambda_m,\Xi,M)$ at matched $\norm{w}_{1,b}$, compare realized gap floors and recursion drift.  
\textbf{Moonshine-DRMM:} modular linear ties among $w_p$ (e.g., by residue classes); test variance reduction of realized slopes without hurting certified floors.  
\textbf{Ablation:} remove each budget in turn at fixed $\norm{w}_{1,b}$ and measure whether certification bands become vacuous or realized performance degrades.

% =========================
% 9. Claims
% =========================
\section{What We Claim (and What We Don't)}
We claim certified, finite-$P$ control of targeted eigen-gaps and slopes under explicit operator-norm budgets and DRMM-lawfulness penalties, plus a scaling study to $P\in\{23,97\}$. We do \emph{not} claim asymptotic ($P\to\infty$) behavior or resolution of deep conjectures. The contribution is a falsifiable, reproducible pipeline with ablations demonstrating when DRMM-style constraints provide genuine advantage.

% =========================
% Appendix (optional)
% =========================
\appendix
\section{Mini-Glossary}
\begin{description}[leftmargin=2.5em, style=nextline]
\item[Plant / Controller] $X_P$ is the system; $C(\omega;w)$ is the designed control.
\item[Bohr--prime multiplier] Rank-1/low-rank channel keyed by a prime-phase vector $u_p(\omega)$.
\item[Woodbury updates] Fast rank-$r$ inverse/solve updates used during verification.
\item[Small-gain] Recursion gain $\eta$ kept strictly within the spectral separation margin.
\item[Davis--Kahan / Weyl] Eigen-projection and eigenvalue perturbation bounds controlling drift and separation.
\end{description}

\section{Certified Spectral Control}
\label{sec:intro}
% P1: Big picture
The interplay between arithmetic structure and spectral phenomena has long been anticipated in number theory and physics. Here we present a minimal, reproducible framework to \emph{control} such spectra with \emph{certificates} using convex optimization.

% P2: Challenge
Directly optimizing spectral quantities depending on full eigendecompositions is nonconvex and costly. We therefore design convex surrogates that bound gaps and slopes from below/above, enabling tractable controller design with rigorous guarantees.

% P3: Contribution
\textbf{Contributions.} (i) A plant/controller framework for prime-scripted operators; (ii) certified convex objectives $J^{(\mathrm{LB})}_{\mathrm{gap}}$, $J^{(\mathrm{UB})}_{\mathrm{slope}}$; (iii) a pilot pinning \\& slope theorem for bounded drives; (iv) a $P{=}5$ case study showing certified gap modulation and null controls.

% P4: Outline
\textbf{Outline.} \cref{sec:model} defines the plant. \cref{sec:theory} proves the pilot theorem. \cref{sec:controller} presents the convex programs. \cref{sec:experiments} shows $P{=}5$ results. \cref{sec:discussion} discusses scope and limitations.

\subsection{Model (Plant): Prime-Scripted Operator}
\label{sec:model}
Let $\PP_P=\{p\in\PP: p\le P\}$ and consider $\mathcal U_P(\omega)=X_P + C(\omega)$ acting on $\ell^2(\PP_P)$.
\begin{definition}[Prime block]
For parameters $\Phi>0$, $\sigma>0$, define
\[
 (X_P)_{pp}=p^{-1/2},\qquad (X_P)_{pq}=\Phi\,\Re\,\zeta\!\Big(\tfrac12+i\log\frac{p}{q}\Big)\,\exp\!\Big(-\tfrac{(\log(p/q))^2}{2\sigma^2}\Big),\quad p\ne q.
\]
\end{definition}
\begin{definition}[Bohr--prime drive]
Given weights $w\in\RR^{|\PP_P|}$, define $m(\omega;w)=\sum_{p\le P}w_p e^{-i\omega\log p}$ and $C(\omega;w)$ by $(\widehat{C f})(\omega)=m(\omega;w)\widehat f(\omega)$. We control $w$ under budgets $\|w\|_1\le W$, $\|w\|_2\le W_2$.
\end{definition}
\begin{remark}[Operator norm]
We use the \emph{operator} norm $\opnorm{\cdot}$ for Weyl/Kato bounds; Hilbert--Schmidt norms enter only in optional factorization/regularization.
\end{remark}

\subsection{Pilot Theorem: Pinning and Slope Bounds}
\label{sec:theory}
Let $S$ be a set of simple eigenvalues of $X_P$ with minimal gap $\delta_S>0$ from the rest of the spectrum.
\begin{theorem}[Pilot pinning \\& slope]
\label{thm:pilot}
If $C\in C^1(\Omega;\mathcal B)$ with $\sup_{\omega\in\Omega}\opnorm{C(\omega)}<\delta_S/2$, then for each $\lambda_i\in S$ there exists a unique $C^1$ branch $\lambda_i(\omega)$ such that
\[
\sup_{\omega\in\Omega}\big|\lambda_i(\omega)-\lambda_i\big|\le \sup_{\omega\in\Omega}\opnorm{C(\omega)},\qquad
\sup_{\omega\in\Omega}\big|\partial_\omega\lambda_i(\omega)\big|\le \sup_{\omega\in\Omega}\opnorm{\partial_\omega C(\omega)}.
\]
Moreover, the Riesz projection onto the band $S$ is well-defined for all $\omega$, and the Davis--Kahan principal angle between invariant subspaces at $0$ and $\omega$ is at most $2\|C\|/\delta_S$.
\end{theorem}
\begin{proof}[Sketch]
Analytic perturbation for simple eigenvalues (Kato) and Weyl inequality yield the bounds; subspace stability follows from Davis--Kahan.
\end{proof}

\section{Controller: Convex Design with Certificates}
\label{sec:controller}
We design $w$ using convex surrogates that bound the realized objectives.
\subsection{Certified lower bound on a target gap}
Let $\Delta_i^{(0)}$ denote gaps of $X_P$ on an index set $\mathcal I$. Define
\[
J_{\mathrm{gap}}^{(\mathrm{LB})}(w):=\min_{i\in\mathcal I}\max\Big\{0,\,\Delta_i^{(0)}-2\,\sup_{\omega\in\Omega}\opnorm{C(\omega;w)}\Big\}.
\]
Since $\sup_\omega\opnorm{C(\omega;w)}\le \|w\|_1$ for the multiplier class, maximizing $J_{\mathrm{gap}}^{(\mathrm{LB})}$ under $\ell_1/\ell_2$ budgets is convex.

\subsection{Certified upper bound on total slope}
\[
J_{\mathrm{slope}}^{(\mathrm{UB})}(w):=\sum_{i\in\mathcal I}\int_{\Omega}\opnorm{\partial_\omega C(\omega;w)}\,d\omega,\quad \text{with}\quad \opnorm{\partial_\omega C}\le \sum_{p\le P}|w_p|\,\log p.
\]
This yields a convex penalty that discourages excessive sensitivity.

\subsection{Convex program}
\begin{align*}
\max_{w\in\RR^{|\PP_P|}}\quad & J_{\mathrm{gap}}^{(\mathrm{LB})}(w) - \lambda\, J_{\mathrm{slope}}^{(\mathrm{UB})}(w)\\
\text{s.t.}\quad & \|w\|_1\le W,\; \|w\|_2\le W_2,\; w\in\mathcal{W}_{\rm phys}.
\end{align*}
Here $\mathcal{W}_{\rm phys}$ encodes platform constraints (bandwidth, symmetry).

\section{Numerical Case Study: $P=5$}
\label{sec:experiments}
\paragraph{Setup.} $\PP_5=\{2,3,5,7,11\}$; choose $(\Phi,\sigma)$; budgets $(W,W_2)$. Target a specific internal gap among the ordered eigenvalues of $X_P$.

\paragraph{Optimization.} Solve the convex program for $w^*$ (projected gradient, coordinate descent). Report $J_{\mathrm{gap}}^{(\mathrm{LB})}(w^*)$ and $J_{\mathrm{slope}}^{(\mathrm{UB})}(w^*)$.

\paragraph{Verification.} Perform one eigentracking sweep to compute the realized gap and compare against the certificate.

\begin{figure}[t]
  \centering
  \includegraphics[width=0.8\textwidth]{fig_spectrum_before_after_placeholder}
  \caption{Eigenvalue trajectories before (baseline) and after optimization ($w^*$); highlighted target gap.}
  \label{fig:gap}
\end{figure}

\begin{figure}[t]
  \centering
  \includegraphics[width=0.7\textwidth]{fig_w_star_placeholder}
  \caption{Optimized prime-lag weights $w^*$ under $\ell_1/\ell_2$ budgets.}
  \label{fig:wstar}
\end{figure}

\begin{table}[t]
  \centering
  \begin{tabular}{lccc}
  \toprule
   & Certified LB & Realized & Null (shuffled lags) \\
  \midrule
  Target gap size & $J_{\mathrm{gap}}^{(\mathrm{LB})}(w^*)$ & $\Delta_{\rm realized}$ & $\Delta_{\rm null}$ \\
  \bottomrule
  \end{tabular}
  \caption{Certificate vs. realized gap and null-control outcome.}
  \label{tab:cert}
\end{table}
\begin{table}[h]
\centering
\caption{Prime vs.\ shuffled delays at $P{=}5$. Shuffled entries are means over 64 seeds.}
\begin{tabular}{ccccccc}
\toprule
$W$ & Prime gap & Shuffled gap (mean) & Gap adv.\% & Prime corr. & Shuffled corr. (mean) & Corr. adv.\% \\
\midrule
0.5 & 1.411 & 1.428 & +1.2 & 0.712 & 0.409 & +74.1 \\
1.0 & 1.328 & 1.349 & +1.5 & 0.712 & 0.344 & +107.3 \\
2.0 & 1.166 & 1.192 & +2.3 & 0.712 & 0.335 & +112.6 \\
4.0 & 0.798 & 0.729 & $-8.6$ & 0.712 & 0.360 & +97.8 \\
\bottomrule
\end{tabular}
\end{table}

\section{Discussion and Outlook}
\label{sec:discussion}
We emphasized methodology: certified convex design and a pilot theorem ensuring spectral stability. Future work adds state-dependent rank-1 mixing, soft lawfulness penalties, and scaling to larger $P$ and phase-diagram studies.

\section*{Reproducibility}
Code and data to be released at: \url{https://example.org/prime-spectral-control}.

\section*{Acknowledgments}
% funding, collaborators

\appendix
\section{Notes on Norm Bounds and Multipliers}
\label{app:norms}
Derive $\sup_\omega\opnorm{C(\omega;w)}\le \|w\|_1$ and $\sup_\omega\opnorm{\partial_\omega C}\le \sum_p |w_p|\log p$ for the convolution multiplier on $L^2$.

\section{Implementation Details}
\label{app:impl}
Optimization uses projected gradient with adaptive step sizes; eigentracking uses Lanczos with warm starts; crossing detection via sign changes in discrete gap derivatives.

% ---------- MATHEMATICAL APPENDIX ----------
\appendix
\section*{Mathematical Appendix: Explicit Proofs and Operator-Norm Bounds}
\addcontentsline{toc}{section}{Mathematical Appendix: Explicit Proofs and Operator-Norm Bounds}

\section{Fourier–Multiplier Bounds for the Temporal Sieve}
\label{app:multiplier}

\begin{lemma}[Operator norm of $C(\omega;w)$]
\label{lem:multiplier-L2}
Let $C(\omega;w)$ act on $L^2(\RR_t)$ as the Fourier multiplier with symbol
$m(\omega;w)=\sum_{p\le P} w_p e^{-i\omega\log p}$. Then for all $\omega$,
\[
\|C(\omega;w)\|_{L^2\to L^2} = \|m(\cdot;w)\|_{L^\infty(\RR)} \le \sum_{p\le P}|w_p| = \|w\|_1.
\]
\end{lemma}

\begin{proof}
$C$ is unitary-conjugate to multiplication by $m$ in the frequency domain:
$\widehat{Cf}(\omega)=m(\omega;w)\widehat f(\omega)$. For multiplication on $L^2$, the operator norm equals the $L^\infty$ norm of the symbol. Since $|e^{-i\omega\log p}|=1$,
$|m(\omega;w)|\le \sum_p |w_p|$ for all $\omega$, proving the bound.
\end{proof}

\begin{lemma}[Derivative bound]
\label{lem:multiplier-derivative}
Assuming $w$ is independent of $\omega$, the weak derivative $\partial_\omega C$ exists as a multiplier with symbol $\partial_\omega m(\omega;w)=-i\sum_{p\le P}w_p(\log p)e^{-i\omega\log p}$ and
\[
\|\partial_\omega C(\omega;w)\|_{L^2\to L^2}
= \|\partial_\omega m(\cdot;w)\|_{L^\infty}
\le \sum_{p\le P}|w_p|\log p.
\]
\end{lemma}

\begin{proof}
Differentiate the symbol termwise; take $L^\infty$ as in \cref{lem:multiplier-L2}.
\end{proof}

\section{Pilot Pinning and Slope: Full Proof}
\label{app:pilot-proof}

We restate \cref{thm:pilot} and give a complete argument.

\begin{theorem*}[Pilot pinning \& slope, restated]
Let $X_P$ be self-adjoint on $\ell^2(\PP_P)$ with a set $S$ of simple eigenvalues $\{\lambda_i\}_{i\in S}$ separated from $\spec(X_P)\setminus S$ by $\delta_S>0$. Let $C\in C^1(\Omega;\mathcal B)$ and
$\sup_{\omega\in\Omega}\|C(\omega)\|<\delta_S/2$. For each $i\in S$ there exists a unique $C^1$ eigenbranch $\lambda_i(\omega)$ of $\mathcal U_P(\omega)=X_P+C(\omega)$ such that
\[
\sup_{\omega\in\Omega}|\lambda_i(\omega)-\lambda_i|
\le \sup_{\omega}\|C(\omega)\|,\qquad
\sup_{\omega\in\Omega}\big|\partial_\omega\lambda_i(\omega)\big|
\le \sup_{\omega}\|\partial_\omega C(\omega)\|.
\]
Moreover the Riesz spectral projection onto $S$ is well-defined for all $\omega$, and the principal angle $\Theta(\omega)$ between the invariant subspaces at $0$ and $\omega$ satisfies
$\sin\Theta(\omega)\le 2\|C\|/\delta_S$.
\end{theorem*}

\begin{proof}
\emph{Existence and uniqueness of branches.} By Kato's analytic perturbation theory for isolated simple eigenvalues, if $\|C(\omega)\|<\delta_S/2$, each $\lambda_i$ extends to a unique $C^1$ branch (indeed analytic in $\omega$ if $C$ is analytic), remaining isolated from the rest of the spectrum.

\emph{Eigenvalue displacement.} For each fixed $\omega$, Weyl's inequality gives
$|\lambda_i(\omega)-\lambda_i|\le \|C(\omega)\|$. Taking $\sup_\omega$ yields the first bound.

\emph{Slope bound.} Differentiate the eigenvalue equation
$(X_P+C(\omega))u_i(\omega)=\lambda_i(\omega)u_i(\omega)$; taking inner product with $u_i(\omega)$ and using $\|u_i(\omega)\|=1$ yields the Hellmann–Feynman identity
$\partial_\omega\lambda_i(\omega)= \langle u_i(\omega),(\partial_\omega C(\omega))u_i(\omega)\rangle$, hence
$|\partial_\omega\lambda_i(\omega)|\le \|\partial_\omega C(\omega)\|$.

\emph{Subspace stability.} Let $E(\omega)=C(\omega)$ and $P_S(0),P_S(\omega)$ be the Riesz projections onto the bands at $0,\omega$. By the Davis–Kahan sin$\Theta$ theorem,
$\| \sin\Theta(\omega)\|\le \|E(\omega)\|/\mathrm{dist}(S,\spec(X_P)\setminus S)$.
Since the minimal gap is $\delta_S$ and we assumed $\sup_\omega\|E(\omega)\|<\delta_S/2$, the claimed bound with factor $2$ follows from standard norm inequalities for projectors.
\end{proof}

\section{Relative Perturbation Determinants for Sums Across Tensor Factors}
\label{app:reldet}

\begin{definition}[Relative perturbation determinant]
Let $A_0$ be a closed operator with nonempty resolvent set; assume $K(z):=(A-A_0)(A_0-z)^{-1}\in\mathcal{S}_2$ (Hilbert–Schmidt) for some $z$. The \emph{2-determinant} is
\[
D(z;A|A_0):=\det\nolimits_{2}\!\big(I+K(z)\big):=\exp\!\Big(\mathrm{tr}\big(\log(I+K(z))-K(z)\big)\Big).
\]
It is analytic in $z$ and encodes the discrete spectral shift of $A$ relative to $A_0$.
\end{definition}

\begin{theorem}[Factorization up to an entire, nonvanishing prefactor]
\label{thm:reldet-factor}
Let $A_0=X_P\otimes\Id\otimes\Id$ on $\ell^2(\PP_P)\otimes L^2(\RR)\otimes\CC^d$. Let
\[
A = A_0 + \underbrace{\Id\otimes C\otimes\Id}_{:=B_1} + \underbrace{\Id\otimes\Id\otimes \Xi}_{:=B_2}.
\]
Assume that for some $z$ in the resolvent of $A_0$,
$K_j(z):=B_j(A_0-z)^{-1}\in\mathcal S_2$ ($j=1,2$).
Then
\[
D(z;A|A_0)
=\mathcal G(z)\,\det\nolimits_2(I+K_1(z))\,\det\nolimits_2(I+K_2(z)),
\]
where $\mathcal G$ is entire and nonvanishing.
\end{theorem}

\begin{proof}[Sketch]
Use the identity $(I+K_1+K_2)=(I+K_1)(I+(I+K_1)^{-1}K_2)$ and the multiplicativity of $\det_2$ up to an analytic, nonzero prefactor that compensates for the $\mathrm{tr}(K)$ terms in the definition of $\det_2$. Since $K_1,K_2\in\mathcal S_2$, the cross-term $(I+K_1)^{-1}K_2\in\mathcal S_2$, and the product formula holds with $\mathcal G(z)=\exp\big(-\mathrm{tr}(K_1K_2)+\cdots\big)$ given by an absolutely convergent series (Plemelj–Smithies). Analyticity and nonvanishing follow from standard arguments in trace-ideal determinant theory.
\end{proof}

\begin{remark}
The naive factorization $\det_2(I-(A_0+B_1+B_2))=\det_2(I-(A_0+B_1))\det_2(I-(A_0+B_2))$ is \emph{false} in general; \cref{thm:reldet-factor} is the correct relative statement.
\end{remark}

\section{Rank-$k$ Updates, Woodbury, and Eigenvalue Perturbations}
\label{app:woodbury}

\begin{lemma}[Woodbury for resolvents]
Let $A=A_0+UV^\ast$ with $U,V:\mathbb{C}^k\to\mathcal H$ and $z\in\rho(A_0)$. Then
\[
(A-z)^{-1}=(A_0-z)^{-1}-(A_0-z)^{-1}U\,(I_k+V^\ast(A_0-z)^{-1}U)^{-1}V^\ast(A_0-z)^{-1}.
\]
\end{lemma}

\begin{proof}
Standard Woodbury identity; verify by multiplication.
\end{proof}

\begin{proposition}[Hoffman–Wielandt/Weyl bounds]
Let $A$ be normal and $E$ be a perturbation. Then
\[
\mathrm{dist}_2(\spec(A+E),\spec(A))\ \le\ \|E\|_{\HS},
\qquad
\mathrm{dist}_\infty(\spec(A+E),\spec(A))\ \le\ \|E\|.
\]
\end{proposition}

\begin{proof}
Classical inequalities (Hoffman–Wielandt for normal matrices; Weyl for general).
\end{proof}

\section{Bounds for the State-Dependent Coupling}
\label{app:state-dep-bounds}

Recall $\mathcal K(\omega)=\Re m(\omega)\,(X_PP_v+P_vX_P)+\Im m(\omega)\,(X_PJ-JX_P)$.

\begin{lemma}[Operator norms]
\label{lem:K-norm}
For any projector $P_v$ and bounded $J$,
\[
\|\mathcal K(\omega)\|
\ \le\ |\Re m(\omega)|\,\|X_PP_v+P_vX_P\| + |\Im m(\omega)|\,\|[X_P,J]\|
\ \le\ 2|\Re m(\omega)|\,\|X_P\| + 2|\Im m(\omega)|\,\|X_P\|\,\|J\|.
\]
\end{lemma}

\begin{proof}
$\|X_PP_v+P_vX_P\|\le \|X_PP_v\|+\|P_vX_P\|\le 2\|X_P\|$ since $\|P_v\|=1$.
For the commutator, $\|[X_P,J]\|\le \|X_PJ\|+\|JX_P\|\le 2\|X_P\|\,\|J\|$.
\end{proof}

\begin{corollary}[Smallness condition]
If $\varepsilon\,\sup_\omega \|\mathcal K(\omega)\|<\delta_S/2$, the band $S$ remains isolated for the active model
$\mathcal U_P(\omega)=X_P+C(\omega)+\varepsilon\mathcal K(\omega)$ and the pilot theorem applies with $C\leadsto C+\varepsilon\mathcal K$.
\end{corollary}

\section{Validity of Certified Surrogates}
\label{app:certificates}

\begin{proposition}[Gap lower bound surrogate]
\label{prop:gapLB}
Let $\Delta_i^{(0)}$ be the $i$-th internal gap of $X_P$. For $E(\omega)=C(\omega;w)+\varepsilon\mathcal K(\omega)$,
\[
\Delta_i(\omega)\ \ge\ \Delta_i^{(0)} - 2\|E(\omega)\|\quad\Rightarrow\quad
\min_{\omega\in\Omega}\Delta_i(\omega)\ \ge\
\Delta_i^{(0)} - 2\sup_{\omega\in\Omega}\|E(\omega)\|.
\]
Thus $J_{\mathrm{gap}}^{(\mathrm{LB})}$ is a rigorous certificate.
\end{proposition}

\begin{proof}
Apply Weyl’s inequality to the two neighboring eigenvalues defining the gap.
\end{proof}

\begin{proposition}[Slope upper bound surrogate]
\label{prop:slopeUB}
For any differentiable branch $\lambda_i(\omega)$,
$|\partial_\omega\lambda_i(\omega)|\le \|\partial_\omega E(\omega)\|$. Hence
\[
\sum_{i\in\mathcal I}\int_\Omega |\partial_\omega\lambda_i(\omega)|\,d\omega
\ \le\ \sum_{i\in\mathcal I}\int_\Omega \|\partial_\omega E(\omega)\|\,d\omega,
\]
justifying $J_{\mathrm{slope}}^{(\mathrm{UB})}$.
\end{proposition}

\begin{proof}
Hellmann–Feynman as in the pilot theorem, then integrate and sum.
\end{proof}

\section{Soft Lawfulness and Subspace Stability}
\label{app:lawful-stability}

\begin{proposition}[Davis–Kahan with lawfulness penalty]
Let $\mathcal S_{\rm law}$ be the Riesz subspace for the “lawful band” of $X_P$. Suppose the controller enforces budgets $W,W_2$ and taper $(\Phi,\sigma)$ so that $\sup_\omega\|E(\omega)\|\le \varepsilon_E$ and the minimal band gap inside $X_P$ satisfies $\delta_S>2\varepsilon_E$. Then the invariant subspace for $\mathcal U_P(\omega)$ over the same index set obeys
\[
\sin\Theta(\omega)\ \le\ \frac{\varepsilon_E}{\delta_S-\varepsilon_E}\ \le\ \frac{2\varepsilon_E}{\delta_S}.
\]
\end{proposition}

\begin{proof}
Davis–Kahan in the version for two self-adjoint operators with separated spectra (triangle inequality on gaps).
\end{proof}

\section{Two-Level Avoided Crossing Model and Arithmetic Fingerprints}
\label{app:two-level}

Consider a pair of modes $(\psi_i,\psi_j)$ near resonance; project $\mathcal U_P(\omega)$ onto $\mathrm{span}\{\psi_i,\psi_j\}$:
\[
H_{\mathrm{eff}}(\omega)=
\begin{pmatrix}
\lambda_i^{(0)}+\delta_i(\omega) & \Gamma_{ij}(\omega)\\
\overline{\Gamma_{ij}(\omega)} & \lambda_j^{(0)}+\delta_j(\omega)
\end{pmatrix},
\quad
\Gamma_{ij}(\omega)\approx
\varepsilon\,\Im m(\omega)\,\langle \psi_i,[X_P,J]\psi_j\rangle
+\langle \psi_i, C(\omega)\psi_j\rangle.
\]
The minimal gap is
\[
\mathrm{gap}_{\min}=2\min_\omega |\Gamma_{ij}(\omega)|\quad\text{when}\quad
\lambda_i^{(0)}+\delta_i(\omega)=\lambda_j^{(0)}+\delta_j(\omega).
\]
If $C$ is diagonal-dominant in the $X_P$-basis, the dominant off-diagonal comes from the commutator term and
\[
\mathrm{gap}_{\min}\ \approx\ 2\,|\Im m(\omega_{\min})|\,
\left|\sum_{p\neq q}\overline{\psi_i(p)}\,\psi_j(q)\,\big(X_{pq}J_{qq}-J_{pp}X_{pq}\big)\right|,
\]
which is \emph{sparse} in $(p,q)$ due to the structure of $X_{pq}$.
Since $X_{pq}\propto \Re\zeta(1/2+i\log(p/q))\,e^{-(\log(p/q))^2/(2\sigma^2)}$, the leading channels $(p,q)$ imprint an \emph{arithmetic fingerprint} on the crossing.

\section{Complexity Summary for the Plant–Controller Loop}
\label{app:complexity}

\begin{itemize}[leftmargin=2em]
\item \textbf{Controller step (convex).} Projection and gradient for $w\in\RR^{|\PP_P|}$ under $\ell_1/\ell_2$ budgets: $O(P)$ per iteration; evaluation of $\|C\|$ and $\|\partial_\omega C\|$ uses \cref{lem:multiplier-L2,lem:multiplier-derivative}.
\item \textbf{Verification step.} Lanczos eigentracking across a grid $|\Omega|$: $O(|\Omega|\,P^2)$ matvecs with warm starts; with rank-$k$ update, Woodbury reduces per-$\omega$ work by a low-rank solve $O(k^3)$ plus $O(kP)$ matvecs.
\item \textbf{Certificates.} $J_{\mathrm{gap}}^{(\mathrm{LB})}$ and $J_{\mathrm{slope}}^{(\mathrm{UB})}$ are computed without diagonalization; the realized spectrum is computed once per design.
\end{itemize}

\section{Notation and Trace-Ideal Conventions}
\label{app:notation}
$\|\cdot\|$ denotes the operator norm unless explicitly marked $\|\cdot\|_{\HS}$.
$\mathcal S_2$ is the Hilbert–Schmidt class; $\det_2$ is the Carleman–Fredholm determinant. Davis–Kahan bounds are stated for orthogonal projectors with the principal-angle norm.

% ---------- end of appendix ----------

\nocite{*}
\bibliographystyle{unsrt}
\bibliography{references}
\end{document}
